\subsection{Desempenho global e ganho incremental}
A Tabela~\ref{tab:metricas} reúne as previsões externas das mesmas 77.535 faixas. Todas as configurações Ridge superaram as referências de média e mediana do treino. Entre as configurações, o Combinado apresentou o menor MAE (8,1151), seguido de Profundo (8,1285), Tradicional (8,2957) e Baseline (8,3525). RMSE e $R^2$ mantiveram a mesma ordenação. O resultado mostra redução de erro na avaliação realizada, sem atribuir às letras causalidade sobre o escore.

\begin{table}[htbp]
\centering\small
\caption{Desempenho nas previsões externas reunidas ($n=77.535$).}
\label{tab:metricas}
\begin{tabular}{lrrrr}
\toprule
Configuração & MAE & RMSE & $R^2$ & MAE ponderado \\
\midrule
Baseline & 8,3525 & 10,8255 & 0,5087 & 8,2194 \\
Tradicional & 8,2957 & 10,7538 & 0,5152 & 8,1652 \\
Profundo & 8,1285 & 10,5540 & 0,5330 & 8,0032 \\
Combinado & 8,1151 & 10,5381 & 0,5344 & 7,9908 \\
Média do treino & 12,5621 & 15,4463 & -0,0003 & 12,4181 \\
Mediana do treino & 12,4786 & 15,5708 & -0,0165 & 12,3213 \\
\bottomrule
\end{tabular}
\fonte{Elaboração própria a partir dos agregados auditados do estudo.}
\nota{Pesos iguais no ajuste e na avaliação principal. Última coluna: sensibilidade com pesos originais nas mesmas previsões; não houve retreino ponderado.}
\end{table}


O acréscimo dos 21 descritores ao Baseline reduziu o MAE em 0,0568 ponto (0,68\%). O acréscimo dos embeddings reduziu-o em 0,2240 ponto (2,68\%). Reunindo ambas as representações, a redução frente ao Baseline foi 0,2374 ponto (2,84\%). Assim, a maior parcela do ganho conjunto já aparece na configuração com embeddings. Adicionar os descritores ao Profundo produziu redução de apenas 0,0134 ponto (0,16\%). A dimensão do ganho importa além de sua direção: a melhora média é modesta diante de um erro absoluto ainda próximo de oito pontos.

A ponderação pelos pesos originais de amostragem, aplicada somente às mesmas previsões, preservou a ordenação. O MAE ponderado passou de 8,2194 no Baseline a 7,9908 no Combinado, redução de 2,78\%. Essa concordância indica que o sentido da comparação não depende apenas da contribuição relativa dos estratos na avaliação sem pesos. Não demonstra que os pesos corrijam a seleção de letras ou que o desempenho ponderado se generalize ao catálogo.

O Combinado superou o Baseline nos cinco testes externos; as diferenças de MAE variaram de $-0,2724$ a $-0,2070$. Retirar um componente de cada vez apenas da avaliação, mantendo os modelos fixos, não inverteu o sinal dos contrastes examinados. Esses resultados descrevem estabilidade nos casos e divisões existentes, mas não são cinco réplicas independentes, validação com reajuste a cada retirada ou intervalos de confiança.

Em seis dos vinte ajustes externos Ridge, a penalização selecionada foi o menor valor da grade, $\alpha=0,1$. Isso não demonstra erro de otimização ou que uma grade ampliada necessariamente melhoraria desempenho. A comparação se refere ao procedimento especificado; modificar a grade depois de observar seus testes externos criaria outra etapa de desenvolvimento. Os resultados atuais foram preservados.

\subsection{Incerteza dos contrastes e alcance da evidência}\label{subsec:incerteza}
As diferenças calculadas nas previsões existentes são conhecidas exatamente para esses casos. Quantificar sua variação em avaliações hipotéticas exige um alvo e hipóteses adicionais. A análise auxiliar do Apêndice~\ref{ap:incerteza} utiliza os componentes para permitir dependência interna e condiciona separadamente no par de modelos de cada fold. Os limites da Tabela~\ref{tab:incerteza} têm nível nominal de 95\% nesse modelo auxiliar, com ajuste Bonferroni; não representam a incerteza integral de seleção, treinamento e construção de dados.

\begin{table}[htbp]
\centering\small
\caption{Limites auxiliares condicionais para os contrastes de interesse.}
\label{tab:incerteza}
\begin{tabular}{lrrrr}
\toprule
Contraste & Avaliação & Delta & Inferior & Superior \\
\midrule
C - B & Iguais & -0,23736 & -0,31400 & -0,16073 \\
C - P & Iguais & -0,01340 & -0,03027 & 0,00347 \\
C - B & Originais & -0,22861 & -0,30454 & -0,15268 \\
C - P & Originais & -0,01234 & -0,02919 & 0,00450 \\
\bottomrule
\end{tabular}
\fonte{Elaboração própria a partir dos agregados auditados do estudo.}
\nota{C: Combinado; B: Baseline; P: Profundo. 95\% nominais sob modelo auxiliar de componentes, com Bonferroni por família de dez marginais. Não são ICs da pipeline integral.}
\end{table}


Para Combinado--Baseline, os limites mantêm valores negativos, entre $-0,3140$ e $-0,1607$ ponto. Para Combinado--Profundo, abrangem zero, de $-0,0303$ a $0,0035$. A sensibilidade ponderada leva à mesma interpretação. O ganho adicional dos descritores além dos embeddings permanece inconclusivo nessa análise; isso não demonstra equivalência entre configurações. Também não se declara significância integral do ganho conjunto, pois o alvo condicional não abrange toda a variação do procedimento.

Essa distinção é compatível com a discussão de \citeonline{bates2024}: erro de um modelo ajustado, desempenho médio de um algoritmo e variabilidade da CV não são a mesma quantidade. Um bootstrap das perdas por componente com previsões congeladas variaria a composição da avaliação, mas não repetiria a escolha de hiperparâmetros nem o treinamento. A grande quantidade de componentes, por si só, não resolve a sobreposição entre treinos ou transforma grupos operacionais em uma amostra independente de composições.

\subsection{Popularidade elevada e subestimação}
A Figura~\ref{fig:bandaserro} mostra que a média global encobre erros diferentes ao longo do alvo. Nas 128 faixas aceitas com popularidade de pelo menos 80, o Combinado teve MAE de 29,3711 pontos, muito acima do global. Essas faixas pertencem a 96 componentes e têm popularidade observada entre 80 e 96: 116 estão em 80--84, oito em 85--89 e quatro em 90--100. Portanto, a dificuldade não se restringe a poucos valores próximos de cem.

\figura{erro_popularidade}{MAE do Baseline e do Combinado por banda de popularidade observada.}{fig:bandaserro}{Elaboração própria a partir das previsões externas existentes. As bandas são definidas pelo alvo observado; contagens indicadas na figura.}

Todas as 128 faixas foram subestimadas pelas quatro configurações. No Combinado, a popularidade média observada foi 81,8125, enquanto a previsão média foi 52,4414; a maior previsão nessa cauda foi 77,9977. Definindo o desvio assinado de diagnóstico por $r_i=\widehat y_i-y_i=-e_i$, a mediana foi $-28,7755$, com primeiro e terceiro quartis de $-37,2444$ e $-22,4381$. Houve 101 subestimações de pelo menos vinte pontos e 22 de pelo menos quarenta pontos, evidenciando um padrão amplo na cauda aceita.

Apesar desses erros elevados, a cauda contém apenas 0,1651\% das faixas avaliadas e produz 0,5975\% do erro absoluto total do Combinado. Sua contribuição aditiva ao MAE global é
\begin{equation}\label{eq:contribuicaocauda}
 \frac{1}{77.535}\sum_{i:y_i\ge80}|e_i|=0,04849.
\end{equation}
Para o erro quadrático total, a participação é 1,4591\%. O baixo peso numérico da cauda permite que um bom resultado médio relativo conviva com desempenho fraco nos escores mais altos. Isso limita aplicações voltadas a identificar músicas muito populares.

A compressão das previsões é compatível com regularização e informação insuficiente para distinguir a cauda, mas o diagnóstico não identifica causalmente seu mecanismo. Condicionar nos maiores valores observados pode gerar erros negativos mesmo sem falha de código. Essa análise também difere de estudar calibração condicionada na previsão, $E(Y\mid\widehat Y)$. No domínio inteiro, o Combinado não produziu nenhuma previsão de pelo menos 80.

\FloatBarrier
\subsection{Previsões negativas e restrição ao suporte}
O Combinado produziu 615 previsões negativas, em 317 componentes, com mínimo de $-32,7066$ e mediana de $-1,7873$. Nessas faixas, a popularidade observada tem média 5,3122 e mediana três; 133 respostas são zero e apenas uma é de pelo menos quarenta. A extrapolação se concentra na borda inferior. A regressão Ridge da Equação~\ref{eq:ridge} não impõe restrição de suporte, de modo que valores negativos não constituem, isoladamente, evidência de dados corrompidos.

Como sensibilidade separada, aplicou-se a projeção
\begin{equation}\label{eq:clipping}
 \widetilde y_i=\min\{100,\max\{0,\widehat y_i\}\}.
\end{equation}
Para $y_i\in[0,100]$, ela não aumenta o erro absoluto nem o quadrático de nenhuma observação, pois aproxima uma previsão externa ao intervalo admissível. A melhora decorrente não representa nova informação das letras. A Tabela~\ref{tab:clipping} apresenta essa análise sem substituir o resultado principal.

\begin{table}[htbp]
\centering\small
\caption{Projeção de suporte como sensibilidade pós-resultado.}
\label{tab:clipping}
\begin{tabular}{lrrrr}
\toprule
Configuração & Abaixo de 0 & Acima de 100 & MAE original & MAE projetado \\
\midrule
Baseline & 537 & 0 & 8,35248 & 8,33781 \\
Tradicional & 584 & 1 & 8,29569 & 8,28013 \\
Profundo & 610 & 0 & 8,12852 & 8,11055 \\
Combinado & 615 & 0 & 8,11512 & 8,09701 \\
\bottomrule
\end{tabular}
\fonte{Elaboração própria a partir dos agregados auditados do estudo.}
\nota{Projeção fixa em [0,100], sem substituir as previsões ou métricas principais.}
\end{table}


A ordenação permanece a mesma e o MAE do Combinado passa a 8,0970, redução de 0,0181 ponto. Nenhuma das 128 previsões da cauda é alterada, porque elas já estavam em $[0,100]$. O clipping reduz extrapolações de suporte, mas não resolve a subestimação das faixas de alta popularidade. Como análise posterior, não foi incorporado retroativamente à regra principal de avaliação.

\subsection{Idiomas, heterogeneidade e limitações}
A Tabela~\ref{tab:idiomas} evidencia heterogeneidade linguística. O ganho Combinado--Baseline é favorável nos status inglês e não inglês, mas desfavorável nos 180 casos incertos. Esse status pequeno não permite uma conclusão geral sobre uma língua específica. Os grupos são estimados automaticamente, e nenhuma faixa foi excluída para melhorar a comparação.

\begin{table}[htbp]
\centering\small
\caption{Desempenho por status do idioma estimado.}
\label{tab:idiomas}
\begin{tabular}{lrrrr}
\toprule
Status & n & Baseline & Profundo & Combinado \\
\midrule
Inglês estimado & 56.908 & 7,6840 & 7,5418 & 7,5278 \\
Não inglês estimado & 20.447 & 10,2036 & 9,7491 & 9,7343 \\
Idioma incerto & 180 & 9,4360 & 9,5398 & 9,8564 \\
\bottomrule
\end{tabular}
\fonte{Elaboração própria a partir dos agregados auditados do estudo.}
\nota{Idioma automático não é rótulo humano. Diagnóstico exploratório sem excluir casos.}
\end{table}


Entre os 76 gêneros com métricas publicadas pelos limiares descritivos de pelo menos cem faixas e vinte componentes, 59 favoreceram o Combinado e 17 não. Essas comparações são exploratórias e podem compartilhar componentes entre categorias; não representam uma bateria confirmatória de testes por gênero. O ganho médio não autoriza prometer benefício para todos os idiomas ou estilos.

O alcance das conclusões combina limitações de fonte, integração e avaliação. O catálogo não cobre probabilisticamente toda a plataforma, a data do escore é desconhecida e a disponibilidade de letras é seletiva. As regras de matching preservam ambiguidades, mas não tiveram precisão global estimada por revisão humana probabilística. O agrupamento reduz vazamentos identificados, sem resolver toda identidade musical ou exposição histórica. Por fim, o encoder avaliado é a variante operacional descrita, e não uma garantia abstrata de desempenho de qualquer embedding E5. Esses limites qualificam a evidência obtida sem exigir que resultados de outras populações sejam misturados aos atuais.
